% 系统总览图（第五版，11-22，Yifan 选的"中心大图"方向；底下三块改成和上面同一种卡片，加了不给清单的对照小图和判定比例条）：
%   顶上一条细步骤带（8 步，分 A 参考几何 / B agent 建房间 / C 材质），中间是同一个房间同一视角的三张大图（点云 → 白膜 → 材质版），
%   底下是本文加的三块。标题只写 PIPELINE，不写 existing、不写 AHa-3D（Yifan 的要求）；出处在论文的图注里交代。
% 版式参考了 Codex 生图给的示意，内容全部是真实数据，没有用生成的图。
% 单独编译成 PDF 再放进论文：cd paper/smallobj/pipeline && ../../../.runtime/tex/tectonic system.tex
% 小图由 scripts/paper/pipeline_assets.py 生成（测试房间 42898849）。
\documentclass[border=1.5pt]{standalone}
\usepackage[T1]{fontenc}
\usepackage[scaled=0.95]{helvet}
\renewcommand{\familydefault}{\sfdefault}
\usepackage{amssymb}
\usepackage{graphicx}
\usepackage{xcolor}
\usepackage{fontawesome5}
\usepackage{tikz}
\usetikzlibrary{arrows.meta,positioning,calc,backgrounds,shadows.blur}
\graphicspath{{../figures/pipeline/}}
\definecolor{ink}{HTML}{14213D}
\definecolor{ink2}{HTML}{5B6170}
\definecolor{hair}{HTML}{D9D5CA}
\definecolor{cream}{HTML}{F6F1E4}
\definecolor{gold}{HTML}{B5872B}
\definecolor{blueM}{HTML}{0F4D92}
\definecolor{blueS}{HTML}{3775BA}
\definecolor{blueF}{HTML}{EAF2FB}
\definecolor{greenM}{HTML}{2E7D32}
\definecolor{greenS}{HTML}{6DB56C}
\definecolor{greenF}{HTML}{EAF6EA}
\definecolor{purpM}{HTML}{6A3D9A}
\definecolor{purpS}{HTML}{A98BCB}
\definecolor{purpF}{HTML}{F3EDF9}
\newcommand{\fsH}{\fontsize{9}{10.5}\selectfont}
\newcommand{\fsT}{\fontsize{7}{8.2}\selectfont}
\newcommand{\fsB}{\fontsize{5.7}{6.9}\selectfont}
\newcommand{\fsS}{\fontsize{5.1}{6.1}\selectfont}
\newcommand{\fsC}{\fontsize{4.9}{5.8}\selectfont}
% 圆角小图：中心 x、y，宽、高（cm），文件；\pimg 不画边框
\newcommand{\rimg}[5]{\begin{scope}\clip[rounded corners=1.8pt] (#1-#3/2,#2-#4/2) rectangle (#1+#3/2,#2+#4/2);
  \node[inner sep=0pt] at (#1,#2) {\includegraphics[width=#3cm,height=#4cm]{#5}};\end{scope}
  \draw[black!22, line width=0.3pt, rounded corners=1.8pt] (#1-#3/2,#2-#4/2) rectangle (#1+#3/2,#2+#4/2);}
\newcommand{\pimg}[5]{\node[inner sep=0pt] at (#1,#2) {\includegraphics[width=#3cm,height=#4cm]{#5}};}
% 步骤小卡：左、右，编号，标题（两行以内），说明
\newcommand{\tile}[5]{%
  \draw[draw=hair, line width=0.4pt, rounded corners=2.5pt, fill=white, blur shadow={shadow blur steps=4, shadow xshift=0pt, shadow yshift=-0.6pt, shadow blur radius=1.2pt, shadow opacity=13}] (#1,-0.98) rectangle (#2,-3.3);
  \node[anchor=north west, inner sep=0pt] at (#1+0.09,-1.07) {\tikz[baseline=(n.base)]{\node[fill=ink, text=white, rounded corners=1.3pt, font=\fsC\bfseries, inner xsep=2.1pt, inner ysep=1.5pt] (n) {#3};}};
  \node[anchor=north west, inner sep=0pt, font=\fsC\bfseries, text=ink, align=left] at (#1+0.4,-1.075) {#4};
  \node[anchor=north, inner sep=0pt, font=\fsC, text=ink2, align=center] at ({(#1+#2)/2},-2.6) {#5};}
\newcommand{\band}[4]{% 左、右，字母，文字
  \fill[cream, rounded corners=2.2pt] (#1,-0.5) rectangle (#2,-0.86);
  \node[anchor=west, inner sep=0pt] at (#1+0.1,-0.68) {\tikz[baseline=(n.base)]{\node[circle, fill=gold, text=white, font=\fsS\bfseries, inner sep=0.9pt, minimum size=7.4pt] (n) {#3};}};
  \node[anchor=west, inner sep=0pt, font=\fsB, text=ink] at (#1+0.45,-0.685) {#4};}
% 本文加的模块：左、右，颜色前缀（blue/green/purp），编号，标题。和上面的步骤卡同一种白底细边，顶上一条模块色
\newcommand{\module}[5]{%
  \draw[draw=hair, line width=0.4pt, rounded corners=3pt, fill=white, blur shadow={shadow blur steps=4, shadow xshift=0pt, shadow yshift=-0.6pt, shadow blur radius=1.2pt, shadow opacity=13}] (#1,-7.52) rectangle (#2,-10.5);
  \begin{scope}\clip[rounded corners=3pt] (#1,-7.52) rectangle (#2,-10.5); \fill[#3M] (#1,-7.52) rectangle (#2,-7.58);\end{scope}
  \node[anchor=west, inner sep=0pt] at (#1+0.14,-7.83) {\tikz[baseline=(n.base)]{\node[fill=#3M, text=white, rounded corners=1.3pt, font=\fsB\bfseries, inner xsep=2.5pt, inner ysep=1.6pt] (n) {#4};}};
  \node[anchor=west, inner sep=0pt, font=\fsB\bfseries, text=ink] at (#1+0.5,-7.835) {#5};}
\newcommand{\bul}[3]{\node[anchor=west, inner sep=0pt, font=\fsB, text=ink2] at (#1,#2) {\textcolor{black!35}{\raisebox{0.5pt}{\scalebox{0.65}{$\bullet$}}}\hspace{2.6pt}#3};}
\begin{document}
\begin{tikzpicture}[x=1cm,y=1cm,
  >={Triangle[length=4.4pt,width=4.8pt]},
  mini/.style={->, line width=0.7pt, draw=black!45, >={Triangle[length=2.6pt,width=3pt]}},
  dtag/.style={rounded corners=2.6pt, font=\fsS, inner sep=1.4pt, minimum height=0.2cm},
  lbl/.style={rounded corners=4pt, fill=white, line width=0.5pt, font=\fsS\bfseries, inner xsep=3pt, inner ysep=1.6pt, anchor=west},
  hero/.style={rounded corners=5pt, fill=cream, font=\fsB\bfseries, text=ink, inner xsep=6pt, inner ysep=2.2pt},
  badge/.style={circle, fill=#1, text=white, font=\fsT\bfseries, inner sep=0.8pt, minimum size=9.6pt},
]
% ================= 标题 =================
\fill[gold] (0,0.02) rectangle (0.065,-0.32);
\node[anchor=west, inner sep=0pt, font=\fsH\bfseries, text=ink] (hA) at (0.17,-0.15) {PIPELINE};
\draw[hair, line width=0.5pt] ($(hA.east)+(0.16,0.13)$) -- ($(hA.east)+(0.16,-0.13)$);
\node[anchor=west, inner sep=0pt, font=\fsT, text=ink2] (hA2) at ($(hA.east)+(0.3,-0.01)$) {one handheld video $\rightarrow$ one editable Blender room};
\draw[hair, line width=0.4pt] ($(hA2.east)+(0.15,0.01)$) -- (14,-0.15);

% ================= 步骤带 =================
\band{0}{6.79}{A}{\textbf{REFERENCE GEOMETRY}\hspace{5pt}{\fsS\textcolor{ink2}{computed once per video, no agent}}}
\band{6.94}{12.0}{B}{\textbf{AGENT BUILDS IN BLENDER}}
\band{12.15}{14}{C}{\textbf{APPEARANCE}}

\tile{0}{1.6}{1}{INPUT\\VIDEO}{handheld RGB\\frames}
\foreach \ang/\dx/\f in {9/0.17/frame_2.jpg, 0/0/frame_1.jpg, -9/-0.17/frame_0.jpg}{
  \begin{scope}[shift={(0.8+\dx,-2.02)}, rotate=\ang]
    \fill[white, rounded corners=1.2pt, blur shadow={shadow blur steps=3, shadow yshift=-0.4pt, shadow xshift=0pt, shadow blur radius=0.8pt, shadow opacity=22}] (-0.36,-0.47) rectangle (0.36,0.47);
    \node[inner sep=0pt, rotate=\ang] at (0,0) {\includegraphics[width=0.66cm]{\f}};
  \end{scope}}

\tile{1.73}{3.33}{2}{PI3X\\GEOMETRY}{camera poses,\\dense points}
\pimg{2.53}{-2.02}{1.38}{1.0}{t_points.jpg}

\tile{3.46}{5.06}{3}{SAM 3\\MASKS}{floor, walls,\\glass, mirrors}
\rimg{3.915}{-2.02}{0.66}{0.93}{mask_wall.jpg}
\rimg{4.605}{-2.02}{0.66}{0.93}{mask_floor.jpg}

\tile{5.19}{6.79}{4}{ROOM\\ALIGNMENT}{gravity-aligned,\\metric frame}
\pimg{5.99}{-2.02}{1.38}{1.0}{t_align.jpg}

\tile{6.94}{8.54}{5}{BLENDER\\CODE}{shell, furniture,\\objects}
\begin{scope}
  \clip[rounded corners=1.8pt] (7.05,-1.54) rectangle (8.43,-2.5);
  \fill[ink] (7.05,-1.54) rectangle (8.43,-2.5);
  \node[anchor=north west, inner sep=0pt, font=\fontsize{2.75}{3.45}\selectfont\ttfamily, text=white, align=left] at (7.1,-1.6)
    {\textcolor{gold!55!white}{def} r\_floor\_lamp(p, w, d, h, prm, cls):\\\ \ p.cyl(0, 0, 0, 0.14, 0.03, 'dark')\\\ \ p.cyl(0, 0, 0.03, 0.015, h - 0.33,\\\ \ \ \ \ \ \ \ 'metal', seg=8)\\\ \ \textcolor{gold!55!white}{for} k \textcolor{gold!55!white}{in} range(4):\\\ \ \ \ p.sphere(0, 0, 0.3 + k * 0.25,\\\ \ \ \ \ \ \ \ \ \ \ \ \ 0.03, 'metal', seg=8)\\\ \ p.cyl(0, 0, h - 0.33, w / 2, 0.33,\\\ \ \ \ \ \ \ \ 'accent', r2=w / 2 * 0.62)};
\end{scope}

\tile{8.67}{10.27}{6}{INSPECT,\\REPAIR}{render at cameras,\\overlay, fix}
\rimg{9.125}{-2.02}{0.66}{0.93}{insp_source.jpg}
\rimg{9.815}{-2.02}{0.66}{0.93}{insp_overlay.jpg}

\tile{10.4}{12.0}{7}{MODEL\\REVIEW}{a second model\\checks renders}
\node[rounded corners=3pt, draw=hair, line width=0.35pt, fill=white, font=\fsS, text=ink, inner xsep=4pt, inner ysep=2.4pt, minimum width=1.12cm] at (11.2,-1.8) {\textcolor{greenM}{\faCheckCircle}\hspace{2.5pt}accept};
\node[rounded corners=3pt, draw=hair, line width=0.35pt, fill=white, font=\fsS, text=ink, inner xsep=4pt, inner ysep=2.4pt, minimum width=1.12cm] at (11.2,-2.24) {\textcolor{gold}{\faExclamationCircle}\hspace{2.5pt}findings};

\tile{12.15}{14}{8}{MATERIALS,\\LIGHTING}{geometry\\unchanged}
\rimg{13.075}{-2.02}{1.6}{1.0}{t_mat.jpg}

\foreach \x in {1.665,3.395,5.125,6.865,8.605,10.335,12.075}{\draw[mini] (\x-0.07,-2.02) -- (\x+0.07,-2.02);}
% 5–6–7 是一个回路：有问题就回到写代码
\draw[mini, draw=ink!60, rounded corners=3pt] (11.2,-3.32) -- (11.2,-3.5) -- (7.74,-3.5) -- (7.74,-3.32);
\node[fill=white, inner xsep=2.5pt, inner ysep=0.5pt, font=\fsC, text=ink2] at (9.47,-3.5) {repeat until accepted};

% 三段各自的产出：连线落在下面三张大图的中轴上（x = 2.2、7.0、11.8），横段和回路同一高度
\draw[mini, draw=ink!60, rounded corners=3pt] (3.395,-3.32) -- (3.395,-3.5) -- (2.2,-3.5) -- (2.2,-3.93);
\draw[mini, draw=ink!60, rounded corners=3pt] (7.74,-3.5) -- (7.0,-3.5) -- (7.0,-3.93);
\draw[mini, draw=ink!60, rounded corners=3pt] (13.075,-3.32) -- (13.075,-3.5) -- (11.8,-3.5) -- (11.8,-3.93);

% ================= 中心三张大图：同一个房间、同一个视角 =================
\begin{scope}[yshift=-0.2cm]  % 以下整体下移 2 mm，给上面的连线留位置
\pimg{2.2}{-5.0}{4.25}{2.5}{hero_points.jpg}
\pimg{7.0}{-5.0}{4.25}{2.5}{hero_white.jpg}
\pimg{11.8}{-5.0}{4.25}{2.5}{hero_mat.jpg}
\foreach \x in {4.6,9.4}{
  \fill[gold!38!white] (\x-0.2,-4.92) -- (\x+0.04,-4.92) -- (\x+0.04,-4.8) -- (\x+0.24,-5.0) -- (\x+0.04,-5.2) -- (\x+0.04,-5.08) -- (\x-0.2,-5.08) -- cycle;}
% 白膜大图右下角的小图：同一个房间不给清单时建出来的样子
\draw[draw=blueS, line width=0.45pt, rounded corners=2.5pt, fill=white, blur shadow={shadow blur steps=4, shadow xshift=0pt, shadow yshift=-0.6pt, shadow blur radius=1.2pt, shadow opacity=15}] (8.42,-5.25) rectangle (9.92,-6.39);
\pimg{9.17}{-5.72}{1.38}{0.88}{hero_base.jpg}
\node[inner sep=0pt, font=\fsC, text=ink2] at (9.17,-6.27) {baseline: \textcolor{ink}{\textbf{37}}};
\node[hero] (h1) at (2.2,-6.56) {REFERENCE GEOMETRY};
\node[hero] (h2) at (7.0,-6.56) {EDITABLE WHITEBOX};
\node[hero] (h3) at (11.8,-6.56) {FINISHED SCENE};
\node[anchor=east, inner sep=0pt, font=\fsC, text=ink2] at ($(h2.west)+(-0.1,0)$) {\textcolor{ink}{\textbf{342}} instances};

% ================= 本文加的三块 =================
\fill[gold] (0,-6.96) rectangle (0.065,-7.3);
\node[anchor=west, inner sep=0pt, font=\fsH\bfseries, text=ink] (hB) at (0.17,-7.13) {THIS WORK};
\draw[hair, line width=0.5pt] ($(hB.east)+(0.16,0.13)$) -- ($(hB.east)+(0.16,-0.13)$);
\node[anchor=west, inner sep=0pt, font=\fsT, text=ink2] (hB2) at ($(hB.east)+(0.3,-0.01)$) {three additions};
\draw[hair, line width=0.4pt] ($(hB2.east)+(0.15,0.01)$) -- (14,-7.13);

% --- 1 小物体清单
\module{0}{7.3}{blue}{1}{SMALL-OBJECT CHECKLIST}
\rimg{1.2}{-8.74}{2.0}{1.36}{b1_detect.jpg}
\rimg{3.42}{-8.74}{1.6}{1.26}{b1_boxes.jpg}
\draw[->, line width=0.95pt, draw=blueM] (2.26,-8.74) -- (2.56,-8.74);
\draw[->, line width=0.95pt, draw=blueM] (4.28,-8.74) -- (4.58,-8.74);
\node[anchor=north, align=center, font=\fsB, text=ink2, inner sep=0pt] at (1.2,-9.53) {per-room categories;\\SAM 3: frames + zooms};
\node[anchor=north, align=center, font=\fsB, text=ink2, inner sep=0pt] at (3.44,-9.53) {lifted to 3D, merged\\6 rooms: \textcolor{blueM}{\textbf{3,210 $\rightarrow$ 2,287}}};
\draw[hair, line width=0.35pt, rounded corners=2.5pt] (4.65,-8.03) rectangle (7.16,-9.31);
\foreach \y/\f/\t in {-8.19/c225.jpg/{\node[dtag, fill=blueM, text=white, anchor=west] at (5.08,\y) {\textbf{built}}; \node[anchor=west, font=\fsS, text=ink2, inner sep=0pt] at (5.58,\y) {as \textit{fruit-bowl}};},
                      -8.51/c502.jpg/{\node[dtag, fill=blueS!42!white, text=ink, anchor=west] at (5.08,\y) {duplicate}; \node[anchor=west, font=\fsS, text=ink2, inner sep=0pt] at (5.88,\y) {of \textit{fruit-bowl}};},
                      -8.83/c018.jpg/{\node[dtag, fill=ink2!30!white, text=ink, anchor=west] at (5.08,\y) {part}; \node[anchor=west, font=\fsS, text=ink2, inner sep=0pt] at (5.54,\y) {of \textit{floor-lamp}};},
                      -9.15/c075.jpg/{\node[dtag, fill=white, draw=ink2!45!white, line width=0.3pt, text=ink, anchor=west] at (5.08,\y) {skipped}; \node[anchor=west, font=\fsS, text=ink2, inner sep=0pt] at (5.85,\y) {outside the room};}}{
  \rimg{4.86}{\y}{0.27}{0.27}{\f}
  \t}
% 四种判定的个数，按比例画成一条（六个测试房间合计 2,287 条）：851 建了、647 重复、727 部件、62 跳过
\fill[blueM] (4.65,-9.43) rectangle (5.475,-9.61);
\fill[blueS!42!white] (5.495,-9.43) rectangle (6.117,-9.61);
\fill[ink2!30!white] (6.137,-9.43) rectangle (6.839,-9.61);
\draw[ink2!45!white, line width=0.3pt] (6.859,-9.43) rectangle (6.92,-9.61);
\node[inner sep=0pt, font=\fsC\bfseries, text=white] at (5.062,-9.522) {851};
\node[inner sep=0pt, font=\fsC, text=ink] at (5.806,-9.522) {647};
\node[inner sep=0pt, font=\fsC, text=ink] at (6.488,-9.522) {727};
\node[anchor=west, inner sep=0pt, font=\fsC, text=ink2] at (6.97,-9.522) {62};
\node[anchor=north, align=center, font=\fsB, text=ink2, inner sep=0pt] at (5.905,-9.7) {one decision per entry};
\draw[hair, line width=0.35pt] (0.2,-10.07) -- (7.1,-10.07);
\node[font=\fsB, text=ink2, anchor=west, inner sep=0pt] at (0.2,-10.28) {\textcolor{blueM}{\faRulerCombined}\hspace{3pt}\textbf{Evaluation:} scored offline against every CA-1M box, by object size};

% --- 2 并行构建
\module{7.5}{10.55}{green}{2}{PARALLEL BUILDS}
\foreach \y/\t in {-8.32/desk, -8.74/chair, -9.16/lamp}{
  \node[draw=greenS, fill=greenF, rounded corners=2pt, font=\fsB, text=ink, minimum width=1.0cm, minimum height=0.32cm, inner sep=1pt] (w\t) at (8.3,\y) {\textcolor{greenM}{\fsS\faCube}\hspace{2.5pt}\t};}
\node[draw=greenM, fill=greenF, rounded corners=2pt, font=\fsB, text=ink, align=center, minimum width=0.92cm, minimum height=0.8cm, inner sep=1.5pt] (wr) at (9.82,-8.74) {\textcolor{greenM}{\fsS\faPen}\\single\\[-1pt]writer};
\draw[line width=0.8pt, draw=greenM, rounded corners=2pt] (wdesk.east) -- ++(0.25,0) |- (wr.west);
\draw[line width=0.8pt, draw=greenM, rounded corners=2pt, dashed] (wlamp.east) -- ++(0.25,0) |- (wr.west);
\draw[->, line width=0.8pt, draw=greenM, >={Triangle[length=3.2pt,width=3.6pt]}] (wchair.east) -- (wr.west);
\bul{7.68}{-9.62}{independent objects in parallel;}
\node[anchor=west, inner sep=0pt, font=\fsB, text=ink2] at (7.825,-9.85) {dependents wait}; 
\bul{7.68}{-10.08}{one writer commits}
\node[anchor=west, inner sep=0pt, font=\fsB, text=ink2] at (7.825,-10.31) {reviewed edits};

% --- 3 物理属性
\module{10.75}{14}{purp}{3}{PHYSICAL PROPERTIES}
\rimg{12.375}{-8.74}{2.7}{1.36}{b3_phys.jpg}
\bul{10.93}{-9.62}{VoMP density per object}
\bul{10.93}{-9.9}{rigid-body rest and push tests}
\node[anchor=west, inner sep=0pt, font=\fsB, text=ink2] at (11.075,-10.13) {on the recovered meshes};

% ================= 连接 =================
\draw[->, line width=1.3pt, draw=blueM] (6.35,-7.5) -- (6.35,-6.78);
\node[lbl, draw=blueM, text=blueM, anchor=east] at (6.22,-7.13) {checklist in, decisions out};
\draw[->, line width=1.3pt, draw=greenM] (7.65,-7.5) -- (7.65,-6.78);
\node[lbl, draw=greenM, text=greenM] at (7.78,-7.13) {schedules builds};
\draw[->, line width=1.3pt, draw=purpM] (11.8,-6.78) -- (11.8,-7.5);
\node[lbl, draw=purpM, text=purpM] at (11.93,-7.13) {scene meshes};
\end{scope}
\end{tikzpicture}
\end{document}
